% !TEX program = xelatex
% Five-page result snapshot; previous S7 design PDF remains unchanged.
\documentclass[11pt,a4paper,fontset=none]{ctexart}
\usepackage{hypothesis_generation}
\usepackage{amsmath,tabularx,pdflscape}
\geometry{left=18mm,right=18mm,top=21mm,bottom=20mm,headsep=6mm}
\setmainfont{Helvetica Neue}
\setsansfont{Helvetica Neue}
\setCJKmainfont{PingFang SC}
\setCJKsansfont{PingFang SC}
\setCJKmonofont{PingFang SC}
\definecolor{hypothesis1}{HTML}{0072B2}
\definecolor{evidencecolor}{HTML}{426D85}
\hypersetup{pdftitle={记忆机制实验报告},pdfauthor={本地研究结果记录},urlcolor=hypothesis1}
\fancyhead[L]{\small\color{hypothesis1}记忆机制实验报告}
\fancyhead[R]{\small S6 $\rightarrow$ S7 · 已完成组件实验}
\fancyfoot[C]{\small 2026年9月6日\quad ·\quad \thepage}
\setlength{\parindent}{0pt}
\setlength{\parskip}{5pt}
\linespread{1.14}
\setlist{nosep,topsep=3pt,leftmargin=1.5em}
\newcommand{\subhead}[1]{\vspace{4pt}{\large\bfseries\color{hypothesis1}#1}\par}
\newcommand{\source}[1]{{\small\color{darkgray}依据：#1。}}
% Generated directly from frozen S7_analysis summary; 6 decimals.
\newcommand{\FourSupportTable}{
\begingroup\small\renewcommand{\arraystretch}{1.28}
\begin{tabularx}{\linewidth}{@{}Xrrrr@{}}
\toprule
采样步长 & A0P0 & A0P1 & A1P0 & A1P1\\
\midrule
8（主设置） & 91.907176 & 91.474634 & 90.639705 & 93.271907\\
12（敏感性） & 91.783236 & 93.624146 & 91.603517 & 93.446457\\
\bottomrule
\end{tabularx}
\endgroup}

\newcommand{\EffectTable}{
\begingroup\small\renewcommand{\arraystretch}{1.28}
\begin{tabularx}{\linewidth}{@{}Xrrrr@{}}
\toprule
干预 & 换图(8) & 差值(8，百分点) & 换图(12) & 差值(12，百分点)\\
\midrule
固定A0：只改位置 & 2/8 & -0.432542 & 3/8 & +1.840910\\
固定A1：只改位置 & 4/8 & +2.632202 & 7/8 & +1.842940\\
固定P0：改关联轨迹 & 2/8 & -1.267471 & 6/8 & -0.179719\\
固定P1：改关联轨迹 & 4/8 & +1.797273 & 3/8 & -0.177689\\
原两臂：A0P0$\rightarrow$A1P1 & 4/8 & +1.364731 & 4/8 & +1.663221\\
\bottomrule
\end{tabularx}
\endgroup}

\newcommand{\ReadoutTable}{
\begingroup\small\renewcommand{\arraystretch}{1.28}
\begin{tabularx}{\linewidth}{@{}ll>{\raggedright\arraybackslash}Xrrr@{}}
\toprule
步长 & 地图 & 原候选＋NMS & 原候选无NMS & 全20＋NMS & 全20无NMS\\
\midrule
8 & A0P0 & 91.907176 & 91.499439 & 91.154065 & 93.471823\\
8 & A0P1 & 91.474634 & 91.699835 & 91.154065 & 93.471823\\
8 & A1P0 & 90.639705 & 91.988238 & 91.154065 & 93.471823\\
8 & A1P1 & 93.271907 & 91.988238 & 91.154065 & 93.471823\\
12 & A0P0 & 91.783236 & 91.398219 & 91.154065 & 93.471823\\
12 & A0P1 & 93.624146 & 91.699835 & 91.154065 & 93.471823\\
12 & A1P0 & 91.603517 & 91.988238 & 91.154065 & 93.471823\\
12 & A1P1 & 93.446457 & 91.988238 & 91.154065 & 93.471823\\
\bottomrule
\end{tabularx}
\endgroup}

\newcommand{\GeometryTable}{
\begingroup\small\renewcommand{\arraystretch}{1.28}
\begin{tabularx}{\linewidth}{@{}llXrr@{}}
\toprule
步长 & 地图 & 共同像素MAE(mm) & 中位差均值(mm) & 自身覆盖(\%)\\
\midrule
8 & A0P0 & 495.391820 & 301.974320 & 8.803863\\
8 & A0P1 & 485.124681 & 281.213864 & 8.587091\\
8 & A1P0 & 495.634166 & 302.531311 & 8.425336\\
8 & A1P1 & 476.342057 & 223.209678 & 8.096275\\
12 & A0P0 & 610.098758 & 500.497461 & 3.503649\\
12 & A0P1 & 614.639795 & 530.810299 & 3.245903\\
12 & A1P0 & 610.404333 & 500.497461 & 3.440276\\
12 & A1P1 & 618.652366 & 552.380415 & 3.246315\\
\bottomrule
\end{tabularx}
\endgroup}

\newcommand{\DevelopmentTable}{
\begingroup\small\renewcommand{\arraystretch}{1.28}
\begin{tabularx}{\linewidth}{@{}Xrrrr@{}}
\toprule
B0原读出支持(\%) & A0P0 & A0P1 & A1P0 & A1P1\\
\midrule
8 & 55.894532 & 56.018638 & 55.987231 & 57.136520\\
12 & 55.987231 & 55.987231 & 57.136520 & 57.136520\\
\bottomrule
\end{tabularx}
\endgroup}

\begin{document}

{\LARGE\bfseries\color{hypothesis1}记忆机制实验报告}\par
{\large 同样的平均位置，为什么有时选得更好、有时更差？}\par
\textbf{S7已运行，9827项独立检查通过。}这是同一房间、固定模型输出后的机制诊断；没有生成新视频，也没有证明一种普遍更好的新算法。

\begin{summarybox}[主要发现：结果取决于位置、关联和选图规则怎样组合]
S6把“第一次存下的位置不改”与“随后匹配时取平均”比较，参考支持从91.907\%升到93.272\%。但平均也改变后续照片匹配到哪个地图点，不能只归因于坐标更准。

S7把两件事拆开。主设置中，\textbf{同一平均位置规则，在首写关联上让支持下降0.433个百分点，在均值关联上却上升2.632个百分点。}稀疏设置里，有些最终不换图的查询，拆开后出现相反变化并抵消。再分开候选范围与NMS后，也没有“一律关闭NMS更好”的结果。
\end{summarybox}

\subhead{问题与范围：模型挑的是4张旧照片}
复用已有CUT3R 224 linear预测，每块20张历史RGB用于建图，后4张只读查询用于选参考。三块共12张不同查询：B0四张只作开发描述，B1/B2的8张为主测试。三个时间块属于同一房间，不是三个独立场景。

\textbf{参考支持}衡量：所选4张历史照片的实测几何能共同支持多少有效查询点，深度一致阈值为50毫米。所有方法使用同一目标mask和分母；它不评价视频美观、动态一致性或语义。建图、选图只使用预测几何与预测相机姿态；实测深度和参考轨迹只用于封存后的评分。

\subhead{实际时间：北京时间，2026年9月6日}
\begin{tabularx}{\linewidth}{@{}p{38mm}X@{}}
\toprule
协议冻结 & 01:03:17.688；看到S6后设计，S7运行前冻结。\\
S7执行 & 01:03:18.259—01:05:23.097，约124.84秒。\\
选择全部封存 & 01:05:22.818；其后才读取已核验的测量数组评分。\\
独立审查 & 01:09:47.634—01:09:54.016，约6.38秒。\\
\bottomrule
\end{tabularx}
{\small 时间来自运行/审查元数据；不是人工研究工时，也不是硬件性能基准。}

\subhead{审查通过意味着什么}
\textbf{9827 = 9147复算 + 663完整性 + 15元数据 + 2源码顺序。}独立程序重建24张地图、12条关联路径，从保存的渲染阵列重算96组投票和384次读出与支持；检查次数不是样本数。S6另已通过3453项独立审查。

\textbf{边界：}没有重新推理模型或光栅化渲染surfel多边形，也没有再次从原PNG生成测量支持。S7复用并逐值核对了S6及其独立审查的测量数组；查询内部状态原张量未保存。已验证的是保存证据之间的一致性，不是完整VMem端到端再复现。
\source{S7运行元数据、S7独立审查；S6结果与独立审查}

\newpage
{\LARGE\bfseries\color{hypothesis1}1｜四张地图，拆开两条路径}\par
\textbf{A = 关联轨迹：}哪些历史观测归入同一个地图点。\textbf{P = 最终位置规则：}该点最后放在哪里。A0来自首写在线过程，A1来自在线平均过程；P0取出生位置，P1按贡献帧质心等帧平均。

\begin{tabularx}{\linewidth}{@{}p{21mm}p{43mm}X@{}}
\toprule
\textbf{地图} & \textbf{关联与位置} & \textbf{可以隔离什么}\\
\midrule
A0P0 & 首写归属＋出生位置 & S6首写对角臂，逐值复现。\\
A0P1 & 同一A0＋平均位置 & 与A0P0只差最终位置。\\
A1P0 & 均值归属＋出生位置 & 与A1P1只差最终位置。\\
A1P1 & 同一A1＋平均位置 & S6均值对角臂，逐值复现。\\
\bottomrule
\end{tabularx}
\textbf{H2不是只改来源标签。}跨A同时改变点数、观测分区、出生法线/半径/颜色等属性和来源集合。同A内这些属性严格相同。S6没有旧逐观测事件；事件是S7新记录并独立复核的，重放时不重新寻找邻居。

\subhead{原候选＋NMS下的平均参考支持（\%）}
\FourSupportTable
\subhead{H1位置与H2关联：条件效应不能混成总效果}
\EffectTable
{\small 换图数按最终ID集合计；两列均为同8张测试图。采样步长12只作敏感性，不代替预定的步长8主设置。}

\textbf{H1和H2都出现实例。}主设置的位置变化在A0上为负、在A1上为正；差中之差为\textbf{+3.064744个百分点}。稀疏设置该平均交互只有\textbf{+0.002030个百分点}，仍不能把均值接近零解读为每张查询都一样。

\subhead{B2抵消：总效果为零，仍有中间变化}
稀疏设置的B2四查询沿$A0P0\rightarrow A1P0\rightarrow A1P1$：先改关联使支持分别下降0.550135、0.575796、0.294272、0.615224个百分点；再改位置恰好升回同样数值，最终回到原参考选择。主设置B2四臂则都不换图。下一页展示全部查询，未只挑这一组。

\textbf{解释限度：}这是固定实现中的描述性分解，不是随机分配下的自然因果效应。不同分解路径可给出不同条件效应；不能宣布唯一主导因素，也不能把敏感性中较高的A0P1改成新的主结果。
\source{S7结果、独立聚合及全部逐查询记录}

\clearpage
\newgeometry{margin=10mm}
\begin{landscape}
\thispagestyle{empty}
\vspace*{-3mm}
\begin{center}
\includegraphics[width=0.94\linewidth]{figures/paired_mechanism_effects.png}
\end{center}
\vspace{-3mm}
{\small\textbf{全部8张测试查询；两档共用色标。}正值不是视频改善。步长12的B2轨迹下降与位置上升抵消；完整精度和ID见CSV。}
\end{landscape}
\restoregeometry

\newpage
{\LARGE\bfseries\color{hypothesis1}2｜同样只选4张，规则仍会改变结果}\par
\textbf{H3检查两件事：候选从哪里来，以及怎样排除相近视角。}NMS优先避免选择相机姿态太相似的照片。表中“全20”只表示20张历史都成为候选，\textbf{最后仍取4张}，没有增加参考预算。

\begin{tabularx}{\linewidth}{@{}p{37mm}XX@{}}
\toprule
\textbf{读出} & \textbf{候选范围} & \textbf{最终选择}\\
\midrule
原候选＋NMS & 可见投票产生的展开候选 & 原NMS及阈值放宽\\
原候选，无NMS & 与上一行完全相同 & 排序前4个唯一ID\\
全20＋NMS & 0--19历史各一次 & 同样NMS及放宽\\
全20，无NMS & 0--19历史各一次 & 排序前4个唯一ID\\
\bottomrule
\end{tabularx}
\subhead{平均支持（\%）：两档采样、四地图、四种读出完整展示}
\ReadoutTable

\begin{comparisonbox}[为什么不能说“去掉NMS就更好”？]
只去掉NMS，主设置A1P1的支持\textbf{下降1.283670个百分点}，A1P0却\textbf{上升1.348533个百分点}。敏感性中分别为−1.458219和+0.384721。NMS的作用依赖当前地图和候选，不能笼统判为好或坏。

只放开候选但保留NMS，支持为91.154065\%；全20候选且无NMS才为93.471823\%。S6近姿态对照的优势因此不能单独归给取消NMS。
\end{comparisonbox}

\subhead{这一次对照保持了哪些执行细节？}
没有直接关掉官方NMS开关，因为该分支还会强制加入最后一张历史并改阈值。新的辅助实现保留展开候选、重复计数、float32距离排序、并列次序及姿态规范化，带NMS路径在全部96次地图查询中复现官方顺序。

原投票器的来源首项双加等行为也保留；本轮没有边修算法边追求更高分。并列候选包含同一ID的重复副本，不能全说成不同相机的量化并列。

\textbf{一致性检查：}两种全20候选读出跨地图和采样密度严格相同，符合“不依赖地图”的定义。它们也不同于使用全部20张支持并集得到的96.773162\%代理上界，后者不是4图预算的对照。
\source{S7独立决策与支持复算；完整384项CSV}

\newpage
{\LARGE\bfseries\color{hypothesis1}3｜几何误差下降，参考支持也未必上升}\par
几何比较只使用\textbf{四张地图都覆盖的有效测量像素}。主设置每查询110--218像素，平均覆盖\textbf{1.935790\%}有效目标；稀疏设置38--82像素，仅\textbf{0.532697\%}。这些范围比S6的两臂共同区域更小，不能直接串接S6与S7误差表排名。

\GeometryTable
{\small MAE为平均绝对深度差。“中位差均值”先逐查询取中位，再对8张查询取均值；不是把所有像素混在一起取中位。自身覆盖是各地图覆盖的有效目标比例。}

\textbf{两个方向相反的实例：}主设置A0内取平均后，共同像素MAE从495.392降至485.125毫米，但参考支持下降0.432542个百分点；稀疏设置A0内MAE从610.099升至614.640毫米，支持却上升1.840910个百分点。

同一投影格不保证同一物理表面，遮挡、位姿与稀疏采样都会影响残差。误差数字必须与共同区域和自身覆盖一起读；未计分区域不能默认为正确。

\subhead{开发数据和负结果都保留}
\DevelopmentTable
{\small B0只有4张开发查询，不并入8张测试均值，也未用于本轮调参。随附CSV保留全部384项读出、支持和所选ID；下降、零变化与开发条件均在其中。}

\subhead{文献怎样限制本轮贡献的说法？}
新综述区分“更新几何”“检索参考”“评价/生成”三层：VMem已有来源索引、可见投票与NMS；AnchorWeave已有局部几何记忆、覆盖检索和多锚点条件。\href{https://arxiv.org/html/2506.18903v3}{[VMem，2025]}\href{https://arxiv.org/html/2602.14941v1}{[AnchorWeave，2026]}\par
经典融合也已有带权平均；近期状态更新与本轮显式地图干预不是同一对象。\href{https://thomaswhelan.ie/Whelan15rss.pdf}{[ElasticFusion，2015]} 本轮价值是\textbf{可复核地暴露条件依赖与抵消}；没有由此证明文献空白、普遍方法优越性或视频质量改善。

\textbf{下一步：}先在新的独立场景中预先固定输入与比较，验证这些现象是否重复；需要实际生成视频后，才讨论生成一致性。旧运行前设计PDF保持为阶段快照，本报告不追改原协议或历史结果。

{\small\textbf{查证入口：}项目\texttt{docs/S7\_RESULTS.md}、\texttt{S7\_INDEPENDENT\_AUDIT.md}、\texttt{S6\_RESULTS.md}、\texttt{S6\_INDEPENDENT\_AUDIT.md}、\texttt{LITERATURE\_SYNTHESIS\_V2.md}；输入哈希、真实编译命令、返回码及逐页QA随报告交付。本报告只重新排版现有真实数据图，没有用生图工具改变实验图。}
\end{document}
